\documentclass[letterpaper]{article}

\usepackage[utf8]{inputenc}
\usepackage[T1]{fontenc}
\usepackage{amsmath,amssymb}
\usepackage{booktabs}
\usepackage{graphicx}
\usepackage{hyperref}
\usepackage{xcolor}
\usepackage{multirow}
\usepackage[margin=1in]{geometry}
\usepackage{natbib}
\usepackage{caption}

\title{Where Do Actions Come From? A Representational Decomposition\\of Vision-Language-Action Models}

\author{
Anonymous Authors
}

\date{}

\begin{document}
\maketitle

%==============================================================================
\begin{abstract}
%==============================================================================
Vision-Language-Action models (VLAs) fine-tune large language models on robot demonstration data to produce action predictions, yet the representational origins of their action-relevant structure remain unclear.
We present a systematic representational and functional analysis of OpenVLA, decomposing contributions from the frozen vision encoder, projector, and LLM backbone.
Using Representational Similarity Analysis (RSA), Centered Kernel Alignment (CKA), and functional action probes across two LIBERO benchmarks and three vision encoders, we establish three findings.
\textbf{(1)}~VL-pretrained encoders (SigLIP, CLIP) dominate action-relevant geometric structure (RSA $\rho{=}0.615$ and $0.553$, respectively) far above both the full VLA ($\rho{=}0.501$) and self-supervised DINO ($\rho{=}0.388$), while all encoders carry comparable action information (probe $R^2{\approx}0.70$).
\textbf{(2)}~The LLM degrades distributed image-token geometric alignment (RSA $0.615{\to}0.501$) but learns to concentrate action signals into a readout token ($R^2{=}0.731$); this concentration is training-dependent (untrained readout $R^2{=}0.686 < $ projector $0.706$).
\textbf{(3)}~Fine-tuning does not alter distributed representation structure (trained RSA $0.501$ $\approx$ untrained $0.485$, $p{=}0.20$) but enables readout routing ($\Delta R^2{=}0.045$, $p{<}0.0001$); language instructions affect neither.
All findings replicate on LIBERO-Spatial with cross-benchmark ratios of $0.95$--$0.97$.
These results reframe VLA fine-tuning as learning to \emph{route} pre-existing visual knowledge rather than \emph{creating} new action representations.
\end{abstract}

%==============================================================================
\section{Introduction}
\label{sec:intro}
%==============================================================================

Vision-Language-Action models~(VLAs) have emerged as a promising paradigm for generalist robot manipulation policies.
By fine-tuning a vision-language model (VLM) on robot demonstration data, VLAs such as OpenVLA~\citep{kim2024openvla} achieve impressive zero-shot and few-shot generalization across manipulation tasks.
The prevailing assumption is that fine-tuning injects new action-relevant knowledge into the LLM backbone, enabling it to ground language instructions in motor behavior.
Yet no prior work has systematically measured \emph{where} action-relevant representational structure originates within the VLA pipeline---at the frozen vision encoder, the learned projector, or the fine-tuned LLM.

This question has practical consequences.
If the frozen vision encoder already provides the dominant action-relevant structure, then efforts to improve VLAs should prioritize vision encoder selection and pretraining rather than scaling LLM fine-tuning.
Conversely, if the LLM creates new representations during fine-tuning, standard scaling recipes remain the right lever.

We present a systematic representational decomposition of OpenVLA, measuring the contribution of each component to action-relevant structure using three complementary tools: \textbf{(i)}~Representational Similarity Analysis (RSA), which quantifies geometric alignment between representation distances and action distances; \textbf{(ii)}~Centered Kernel Alignment (CKA), which captures global structure preservation; and \textbf{(iii)}~functional action probes, which measure how much action information is linearly decodable.
We compare trained and untrained VLA weights, with and without language instructions, across two LIBERO benchmarks (Goal, Spatial) and three standalone vision encoders (SigLIP, CLIP ViT-L, DINOv2).

Our analysis yields three core findings:
\begin{enumerate}
    \item \textbf{VL-pretrained encoder dominance.} The frozen SigLIP vision encoder paired with a learned projector produces representations whose geometric structure is most aligned with action structure (RSA $\rho{=}0.615$). An independent VL-pretrained encoder, CLIP ViT-L, achieves a statistically indistinguishable $\rho{=}0.553$ ($p{=}0.08$), while self-supervised DINOv2 falls significantly behind ($\rho{=}0.388$, $p{=}0.001$). This gap is geometric, not informational: all three encoders carry comparable action information ($R^2{\approx}0.70$).
    \item \textbf{LLM redistributes, training-dependently.} The LLM degrades the distributed image-token geometric alignment inherited from the projector (RSA drops from $0.615$ to $0.501$), but trained LLM layers learn to concentrate action signals into a readout token ($R^2{=}0.731$ vs.\ projector $0.706$, $p{=}0.001$). Untrained LLM weights cannot achieve this concentration (readout $R^2{=}0.686 < $ projector $0.706$).
    \item \textbf{Training routes, not creates.} Fine-tuning leaves the distributed image-token representation structure unchanged (trained RSA $0.501$ vs.\ untrained $0.485$, bootstrap $p{=}0.20$) while dramatically improving readout concentration ($\Delta R^2{=}0.045$, $p{<}0.0001$). Language instruction content has no measurable effect on either.
\end{enumerate}

These findings are corroborated by a CKA--RSA dissociation (CKA ordering reverses the RSA ordering), a CLIP-ViT control experiment that generalizes Claim~1 beyond SigLIP, and full replication on LIBERO-Spatial.

%==============================================================================
\section{Related Work}
\label{sec:related}
%==============================================================================

\paragraph{Vision-Language-Action Models.}
OpenVLA~\citep{kim2024openvla} builds on the Prismatic VLM architecture~\citep{karamcheti2024prismatic}, using a SigLIP vision encoder with a learned projector feeding into an LLM backbone fine-tuned on robot demonstrations.
RT-2~\citep{brohan2023rt2} and related systems established that VLMs can be adapted for action prediction, but the representational mechanisms underlying this adaptation remain opaque.

\paragraph{Behavioral Analysis of VLAs.}
Grant et al.~(2026) demonstrate at the behavioral level that the visual pathway dominates VLA action generation.
Hou~(2026, LangGap) and Fei et al.~(2025, LIBERO-Plus) independently observe that VLAs largely ignore language instructions.
Our work provides representational-level explanations for these behavioral findings.

\paragraph{Mechanistic Interpretability of VLAs.}
Haon et al.~(2025) apply activation steering to OpenVLA, finding sparse semantic directions that influence actions.
Lu et al.~(2025) show that linear probes on OpenVLA hidden layers achieve $>$90\% task classification accuracy.
Swann et al.~(2026) use sparse autoencoders, while Molinari et al.~(2025) study emergent world representations.
Rizzardo et al.~(2026) examine feature observability and controllability.
Our approach complements these with a quantitative component-level decomposition using RSA, CKA, and functional probes.

\paragraph{Representation Quality in Fine-Tuned VLMs.}
Fang~(2025) shows that fine-tuning erodes VLM reasoning capabilities.
Our findings extend this observation: the LLM step itself---not just fine-tuning---degrades distributed geometric alignment, though fine-tuning enables compensatory readout concentration.


%==============================================================================
\section{Method}
\label{sec:method}
%==============================================================================

\subsection{Model and Components}

We analyze OpenVLA, which comprises three components:
\textbf{(i)}~a frozen SigLIP ViT-SO400M vision encoder;
\textbf{(ii)}~a learned two-layer MLP projector that maps vision tokens to the LLM embedding space; and
\textbf{(iii)}~a Llama-2-7B LLM backbone fine-tuned on robot demonstrations.
We extract representations at the projector output (``projector''), at each of the 32 LLM layers (``LLM $L_k$''), and at the final pre-action token position (``readout'').

To disentangle training effects, we compare: (a)~the published fine-tuned OpenVLA weights (``trained''), and (b)~the pre-fine-tuning Prismatic VLM weights (``untrained''), which share architecture and vision encoder but have not seen robot data.
We further ablate language instruction presence (with instruction vs.\ no instruction).

\subsection{Benchmarks}

We evaluate on two LIBERO suites~\citep{liu2024libero}:
\textbf{LIBERO-Goal} (10 tasks varying in goal state) and \textbf{LIBERO-Spatial} (10 tasks varying in spatial arrangement).
For each benchmark, we collect 50 demonstrations per task (500 total), extracting paired observations $(o_t, a_t)$ where $o_t$ is the visual observation and $a_t$ is the 7-DoF action.

\subsection{Standalone Vision Encoders}

To test whether action-relevant structure is specific to SigLIP or general to VL-pretrained encoders, we extract features from two additional encoders applied to the same observations:
\textbf{CLIP ViT-L/14}~\citep{radford2021clip} (VL contrastive, softmax loss) and
\textbf{DINOv2 ViT-L/14}~\citep{oquab2024dinov2} (self-supervised, no language).

\subsection{Representational Similarity Analysis (RSA)}

RSA quantifies how well the geometry of a representation space aligns with the geometry of the action space.
For each model condition, we compute the pairwise representational dissimilarity matrix (RDM) over all observation pairs and the corresponding action dissimilarity matrix.
The RSA score is the Spearman rank correlation $\rho$ between the upper triangles of these two matrices.
Higher RSA indicates that observations leading to similar actions are placed nearby in representation space.

We report pooled RSA (all tasks) and per-task RSA, with 95\% bootstrap confidence intervals (10,000 resamples) and permutation-based $p$-values for pairwise comparisons.

\subsection{Centered Kernel Alignment (CKA)}

CKA~\citep{kornblith2019cka} measures global structural similarity between two representation spaces.
Unlike RSA, which is sensitive to local distance structure, CKA captures whether a linear or kernel transformation can align the two spaces.
We compute linear CKA between each model condition's representation and the action matrix.

\subsection{Functional Action Probes}

We train ridge regression probes to predict 7-DoF actions from representations, reporting cross-validated $R^2$.
We evaluate two targets: \textbf{raw actions} $a_t$ and \textbf{action deltas} $\Delta a_t = a_t - a_{t-1}$.
Two aggregation methods are used for LLM representations:
\textbf{image\_mean} (mean over image token positions, for fair comparison with the projector) and
\textbf{last\_preaction} (the final token before action prediction, capturing readout concentration).

%==============================================================================
\section{Experiments}
\label{sec:experiments}
%==============================================================================

\subsection{Vision Encoder Dominance (Claim 1)}
\label{sec:exp-encoder}

Our first analysis tests whether the frozen vision encoder, rather than the LLM backbone, is the primary source of action-relevant geometric structure.

\paragraph{Setup.}
We compute RSA between action dissimilarities and representation dissimilarities for four conditions on LIBERO-Goal: SigLIP+Projector output, full VLA image\_mean (trained, layer 23), DINOv2, and CLIP ViT-L/14.

\paragraph{Results.}
Table~\ref{tab:encoder-rsa} presents the RSA and functional probe results.
SigLIP+Projector achieves the highest geometric alignment with action structure ($\rho{=}0.615$, 95\% CI $[0.53, 0.70]$).
The full VLA's distributed image tokens show reduced alignment ($\rho{=}0.501$), and DINOv2 is substantially lower ($\rho{=}0.388$).

\begin{table}[t]
\centering
\caption{Vision encoder RSA and functional probe comparison on LIBERO-Goal. RSA $\rho$ measures geometric alignment with action structure; $R^2$ measures linearly decodable action information; $\Delta R^2$ is for action-delta prediction.}
\label{tab:encoder-rsa}
\begin{tabular}{@{}lcccc@{}}
\toprule
Model & Pretraining & RSA $\rho$ & 95\% CI & Action $R^2$ \\
\midrule
SigLIP+Proj        & VL (sigmoid)  & \textbf{0.615} & {[0.53, 0.70]} & 0.706 \\
CLIP ViT-L         & VL (softmax)  & 0.553          & {[0.47, 0.63]} & 0.695 \\
VLA img\_mean L23  & VL+FT         & 0.501          & {[0.42, 0.58]} & 0.683 \\
DINOv2             & Self-sup.     & 0.388          & {[0.30, 0.47]} & 0.712 \\
\bottomrule
\end{tabular}
\end{table}

\paragraph{CLIP-ViT Pretraining Control.}
A critical question is whether SigLIP's high RSA reflects properties specific to its sigmoid loss and the SigLIP-specific projector training, or generalizes to VL-pretrained encoders.
CLIP ViT-L achieves RSA $\rho{=}0.553$ (Table~\ref{tab:clip-bootstrap}), with the CLIP--SigLIP difference not significant ($\Delta{=}{-}0.062$, 95\% CI $[-0.14, 0.02]$, $p{=}0.08$), while the CLIP--DINO gap is highly significant ($\Delta{=}0.165$, 95\% CI $[0.08, 0.25]$, $p{=}0.001$).
This establishes a three-tier structure: \textbf{VL-pretrained} (SigLIP $\approx$ CLIP, $\rho{\sim}0.55$--$0.62$) $\gg$ \textbf{VLA backbone} ($\rho{\sim}0.50$) $\gg$ \textbf{self-supervised} (DINO, $\rho{\sim}0.39$).

\begin{table}[t]
\centering
\caption{Pairwise bootstrap tests on RSA (LIBERO-Goal). The CLIP vs.\ SigLIP gap is not significant, confirming that VL pretraining---not SigLIP specifically---drives geometric alignment.}
\label{tab:clip-bootstrap}
\begin{tabular}{@{}lcccc@{}}
\toprule
Comparison & $\Delta$ & 95\% CI & $p$ & Significant \\
\midrule
CLIP vs.\ SigLIP       & $-$0.062 & {[$-$0.14, 0.02]} & 0.08  & No  \\
CLIP vs.\ DINO         & +0.165   & {[0.08, 0.25]}    & 0.001 & Yes \\
CLIP vs.\ VLA img\_mean & +0.052  & {[$-$0.03, 0.13]} & 0.10  & No  \\
\bottomrule
\end{tabular}
\end{table}

\paragraph{Geometric--Functional Dissociation.}
A striking pattern emerges from the functional probes.
Despite large RSA differences, all encoders achieve comparable action prediction $R^2$: DINO $0.712 \approx$ SigLIP $0.706 \approx$ CLIP $0.695$.
Action information is \emph{present} regardless of pretraining, but \emph{geometric alignment} with action structure requires VL pretraining.
VL pretraining organizes visual features by semantic similarity (matching language-described goals), which naturally aligns with task-conditioned action structure.
DINO organizes by visual similarity---sufficient for a probe to extract action information, but geometrically misaligned.


\subsection{LLM Redistributes Action Structure (Claim 2)}
\label{sec:exp-llm}

\paragraph{Setup.}
We track how representation structure changes from the projector output through 32 LLM layers, measuring RSA and probe $R^2$ at each layer using two aggregation methods: image\_mean and last\_preaction.

\paragraph{Geometric Degradation.}
RSA drops sharply from projector ($0.615$) to LLM Layer~1 ($\sim$0.42, a 32\% drop), with a non-monotonic pattern through subsequent layers, peaking at $0.501$ at Layer~23 for image\_mean.
The LLM transforms the inherited geometric structure, reducing local alignment with action distances.

\paragraph{Readout Concentration.}
While image\_mean probe $R^2$ drops from $0.706$ (projector) to $0.683$ (LLM, $\Delta{=}0.023$, $p{=}0.001$), the last\_preaction token at Layer~19 achieves $R^2{=}0.731$---\emph{exceeding} the projector ($p{=}0.001$).
The LLM self-attention mechanism concentrates action-relevant information from distributed image tokens into the readout position, compensating for the geometric degradation.

\paragraph{Training Dependence of Readout.}
This concentration is not an architectural property but a learned behavior.
Untrained LLM weights produce a best last\_preaction $R^2$ of only $0.686$ at Layer~32, below the projector's $0.706$.
The trained--untrained readout gap is substantial ($\Delta{=}0.045$, 95\% CI $[0.032, 0.058]$, $p{<}0.0001$).

\paragraph{CKA--RSA Dissociation.}
CKA reveals a complementary picture: CKA with actions \emph{increases} monotonically through LLM layers (projector $0.081 \to$ Layer~32 $0.098$), and the ordering reverses relative to RSA (CKA: untrained $0.089 >$ trained $0.083 >$ projector $0.081 >$ DINO $0.066$).
This dissociation reflects different aspects of representational transformation: RSA captures local distance alignment (degraded), while CKA captures global structure preservation (maintained or improved).
The highest CKA value occurs at the last\_preaction position ($0.127$), consistent with readout concentration.


\subsection{Training: Routing Without Creating (Claim 3)}
\label{sec:exp-training}

\paragraph{Setup.}
We compare trained and untrained VLA weights across a 2$\times$2 design: \{trained, untrained\} $\times$ \{with instruction, no instruction\}, measuring both RSA (image\_mean) and readout probe $R^2$.

\paragraph{Distributed Representation Invariance.}
Table~\ref{tab:2x2-goal} shows the 2$\times$2 RSA matrix for LIBERO-Goal.
All four conditions fall within a narrow band ($[0.48, 0.50]$).
The training effect on image\_mean RSA is not significant ($\Delta{=}0.016$, bootstrap $p{=}0.20$).
The instruction effect is negligible: $\Delta{=}0.004$ for trained, $0.003$ for untrained.
Random, shuffled, and empty instruction controls confirm this invariance (all within $\pm 0.017$ of the baseline).

\begin{table}[t]
\centering
\caption{2$\times$2 RSA matrix (image\_mean, LIBERO-Goal). Training and instruction effects are both non-significant. All values fall in a narrow band $[0.48, 0.50]$.}
\label{tab:2x2-goal}
\begin{tabular}{@{}lccc@{}}
\toprule
Condition & With Inst. & No Inst. & Instruction $\Delta$ \\
\midrule
Trained   & 0.501 & 0.501 & 0.000 \\
Untrained & 0.485 & 0.482 & 0.003 \\
\midrule
Training $\Delta$ & 0.016 & 0.019 & \\
\bottomrule
\end{tabular}
\end{table}

\paragraph{Readout Routing.}
In contrast, the readout token shows a dramatic training effect.
Trained last\_preaction achieves $R^2{=}0.731$ vs.\ untrained $0.686$ ($\Delta{=}0.045$, $p{<}0.0001$).
Fine-tuning teaches the LLM to route action-relevant information to the prediction position via attention, without modifying the underlying distributed representations.
This is consistent with the next-token prediction training objective: the model learns where to ``look'' in the context, not how to create new features.


\subsection{Cross-Benchmark Validation: LIBERO-Spatial}
\label{sec:exp-spatial}

To assess generalization beyond a single benchmark, we replicate the full analysis pipeline on LIBERO-Spatial (10 tasks, spatial arrangement variations).

\paragraph{RSA Ordering Replicates.}
Table~\ref{tab:cross-benchmark-rsa} shows RSA values across both benchmarks.
The ordering is perfectly preserved: SigLIP ($0.587$) $\gg$ VLA trained ($0.482$) $\approx$ untrained ($0.468$) $\gg$ DINO ($0.371$).
Cross-benchmark ratios range from $0.955$ to $0.965$, indicating that absolute RSA values are slightly lower on LIBERO-Spatial (spatial tasks have more similar scenes, reducing inter-task discriminability) but relative structure is unchanged.

\begin{table}[t]
\centering
\caption{Cross-benchmark RSA comparison. Ordering is preserved; ratios (Spatial/Goal) are stable at $0.955$--$0.965$.}
\label{tab:cross-benchmark-rsa}
\begin{tabular}{@{}lccc@{}}
\toprule
Model & LIBERO-Goal & LIBERO-Spatial & Ratio \\
\midrule
SigLIP+Proj         & 0.615 & 0.587 & 0.955 \\
VLA trained img\_mean  & 0.501 & 0.482 & 0.962 \\
VLA untrained img\_mean & 0.485 & 0.468 & 0.965 \\
DINOv2               & 0.388 & 0.371 & 0.956 \\
\bottomrule
\end{tabular}
\end{table}

\paragraph{Training Invariance Replicates.}
On LIBERO-Spatial, the training effect on RSA is again non-significant ($\Delta{=}0.014$, 95\% CI $[-0.04, 0.06]$, $p{=}0.22$).
The 2$\times$2 matrix (Table~\ref{tab:2x2-spatial}) shows all conditions within $[0.46, 0.49]$, with instruction effects of $|\Delta|{\leq}0.005$.

\begin{table}[t]
\centering
\caption{2$\times$2 RSA matrix (image\_mean, LIBERO-Spatial). Same pattern as LIBERO-Goal: no significant training or instruction effect.}
\label{tab:2x2-spatial}
\begin{tabular}{@{}lccc@{}}
\toprule
Condition & With Inst. & No Inst. & Instruction $\Delta$ \\
\midrule
Trained   & 0.482 & 0.479 & $-$0.003 \\
Untrained & 0.468 & 0.463 & 0.005 \\
\midrule
Training $\Delta$ & 0.014 & 0.016 & \\
\bottomrule
\end{tabular}
\end{table}

\paragraph{Functional Probes Replicate.}
Table~\ref{tab:cross-benchmark-probe} shows that functional probe results are highly consistent across benchmarks, with all cross-benchmark $\Delta$ values within $0.02$.
Readout concentration remains significant on LIBERO-Spatial: projector vs.\ image\_mean $\Delta{=}0.026$ ($p{=}0.002$); trained vs.\ untrained readout $\Delta{=}0.046$ ($p{<}0.0001$).
The geometric--functional dissociation also replicates: DINO action $R^2{=}0.693 \approx$ projector $0.688$ ($p{=}0.28$), but RSA gap $= 0.216$.

\begin{table}[t]
\centering
\caption{Functional probe $R^2$ across benchmarks. All values are consistent (cross-benchmark $|\Delta|<0.021$).}
\label{tab:cross-benchmark-probe}
\begin{tabular}{@{}lccc@{}}
\toprule
Model & LIBERO-Goal $R^2$ & LIBERO-Spatial $R^2$ & $\Delta$ \\
\midrule
SigLIP+Proj          & 0.706 & 0.688 & $-$0.018 \\
Trained img\_mean     & 0.683 & 0.662 & $-$0.021 \\
Trained readout      & 0.731 & 0.715 & $-$0.016 \\
Untrained readout    & 0.686 & 0.669 & $-$0.017 \\
DINOv2               & 0.712 & 0.693 & $-$0.019 \\
\bottomrule
\end{tabular}
\end{table}

\paragraph{CKA Dissociation Replicates.}
CKA ordering on LIBERO-Spatial mirrors LIBERO-Goal: untrained ($0.085$) $>$ trained ($0.079$) $>$ projector ($0.077$) $>$ DINO ($0.062$).
The reversal relative to RSA ordering holds across benchmarks, confirming that this dissociation reflects intrinsic properties of the LLM transformation rather than benchmark-specific artifacts.

\paragraph{Statistical Summary (LIBERO-Spatial).}
Table~\ref{tab:spatial-bootstrap} summarizes bootstrap tests on LIBERO-Spatial, confirming the same pattern of results as LIBERO-Goal.

\begin{table}[t]
\centering
\caption{Bootstrap tests on LIBERO-Spatial. Training effect on RSA is non-significant; readout concentration effects are highly significant.}
\label{tab:spatial-bootstrap}
\begin{tabular}{@{}lccc@{}}
\toprule
Comparison & $\Delta$ & 95\% CI & $p$ \\
\midrule
Proj vs.\ img\_mean $R^2$ (L18) & 0.026 & {[0.012, 0.041]} & 0.002 \\
Trained vs.\ untrained readout $R^2$ & 0.046 & {[0.030, 0.062]} & $<$0.0001 \\
Training effect on RSA & 0.014 & {[$-$0.04, 0.06]} & 0.22 \\
\bottomrule
\end{tabular}
\end{table}

%==============================================================================
\section{Discussion}
\label{sec:discussion}
%==============================================================================

\paragraph{Reframing VLA Fine-Tuning.}
Our findings challenge the implicit assumption that VLA fine-tuning creates new action-relevant representations.
Instead, the frozen vision encoder---specifically, VL-pretrained encoders---provides the geometric structure that aligns with action space, and fine-tuning teaches the LLM to \emph{route} this pre-existing information to the prediction position.
This is consistent with a growing body of evidence that foundation models acquire rich world knowledge during pretraining that is preserved across downstream tasks.

\paragraph{Why VL Pretraining Matters Geometrically.}
The geometric--functional dissociation (DINO matches SigLIP/CLIP in probe $R^2$ but not RSA) reveals that VL pretraining contributes geometric alignment, not information content.
VL contrastive training organizes visual features by semantic similarity (aligning with language descriptions), which naturally matches task-conditioned action structure (actions are defined by language-described goals).
Self-supervised pretraining organizes by visual similarity---sufficient for probes but geometrically misaligned.

\paragraph{SigLIP vs.\ CLIP.}
SigLIP achieves slightly higher RSA than CLIP ($0.615$ vs.\ $0.553$), though the difference is not significant ($p{=}0.08$).
Two factors may contribute: (i)~SigLIP's sigmoid loss better preserves fine-grained distance structure compared to softmax cross-batch normalization; (ii)~the OpenVLA projector was trained specifically for SigLIP's output space.
The key result is that \emph{both} VL-pretrained encoders substantially outperform DINO, confirming that VL pretraining---not SigLIP specifically---is the critical ingredient.

\paragraph{CKA--RSA Dissociation.}
The reversal of CKA and RSA orderings reveals that the LLM transformation has two distinct effects: it degrades local distance alignment (RSA decreases) while preserving or improving global structural correspondence (CKA increases).
This is consistent with the LLM performing a nonlinear remapping that shuffles local neighborhoods while maintaining large-scale structure---a transformation suited to next-token prediction that needs global context more than local metric preservation.

\paragraph{Implications for VLA Design.}
Three practical implications follow:
\textbf{(i)}~Vision encoder pretraining is the primary lever for action-relevant representation quality; VL contrastive pretraining is preferred over self-supervised.
\textbf{(ii)}~For tasks requiring distributed image-level representations (e.g., spatial reasoning), projector output may be preferable to LLM-processed features.
\textbf{(iii)}~The LLM's value lies in learned attention-based routing, not representation creation---architectures that separate routing from representation may be more efficient.

\paragraph{Explaining Language Instruction Insensitivity.}
Our finding that instruction content has no measurable effect on representation structure provides a representational-level explanation for behavioral observations that VLAs ignore language instructions~\citep{hou2026langgap,fei2025liberoplus}.
If instruction tokens do not alter the representation geometry through which actions are predicted, language grounding must operate through a different mechanism---if it operates at all.

\paragraph{Limitations.}
Our analysis is conducted on OpenVLA, a single VLA architecture. While OpenVLA is representative of the dominant ``VLM + fine-tuning'' paradigm, findings may not transfer to architecturally distinct approaches (e.g., diffusion-based action heads, cross-attention fusion).
We study representation structure, not behavioral performance; higher RSA does not guarantee higher task success rate, though it reflects better-organized representations.
All experiments use LIBERO simulation; transfer to real-world manipulation remains to be verified.

%==============================================================================
\section{Conclusion}
\label{sec:conclusion}
%==============================================================================

We presented a systematic representational decomposition of OpenVLA, revealing that action-relevant geometric structure originates in the frozen VL-pretrained vision encoder, is degraded by LLM processing at the image-token level, and is reconcentrated into a readout token through training-dependent attention routing.
Fine-tuning teaches routing, not representation creation, and language instructions affect neither.
All findings replicate across two LIBERO benchmarks and generalize from SigLIP to CLIP, establishing VL-pretrained encoder dominance as a robust phenomenon.
These results suggest that improving VLA performance should prioritize vision encoder quality and pretraining strategy over LLM scaling.

\bibliographystyle{plainnat}

\begin{thebibliography}{20}

\bibitem[Brohan et~al.(2023)]{brohan2023rt2}
A.~Brohan et~al.
\newblock RT-2: Vision-Language-Action Models Transfer Web Knowledge to Robotic Control.
\newblock \emph{CoRL}, 2023.

\bibitem[Fang(2025)]{fang2025erosion}
H.~Fang.
\newblock Fine-Tuning Erodes VLM Reasoning.
\newblock \emph{arXiv preprint}, 2025.

\bibitem[Fei et~al.(2025)]{fei2025liberoplus}
H.~Fei et~al.
\newblock LIBERO-Plus: Benchmarking Language Grounding in VLAs.
\newblock \emph{arXiv preprint}, 2025.

\bibitem[Grant et~al.(2026)]{grant2026visual}
E.~Grant et~al.
\newblock The Visual Pathway Dominates Action Generation in VLAs.
\newblock 2026.

\bibitem[Haon et~al.(2025)]{haon2025steering}
A.~Haon et~al.
\newblock Activation Steering in OpenVLA.
\newblock \emph{CoRL}, 2025.

\bibitem[Hou(2026)]{hou2026langgap}
X.~Hou.
\newblock LangGap: Language Grounding Failures in Vision-Language-Action Models.
\newblock 2026.

\bibitem[Karamcheti et~al.(2024)]{karamcheti2024prismatic}
S.~Karamcheti et~al.
\newblock Prismatic VLMs: Investigating the Design Space of Visually-Conditioned Language Models.
\newblock \emph{ICML}, 2024.

\bibitem[Kim et~al.(2024)]{kim2024openvla}
M.~Kim et~al.
\newblock OpenVLA: An Open-Source Vision-Language-Action Model.
\newblock \emph{arXiv:2406.09246}, 2024.

\bibitem[Kornblith et~al.(2019)]{kornblith2019cka}
S.~Kornblith, M.~Norouzi, H.~Lee, and G.~Hinton.
\newblock Similarity of Neural Network Representations Revisited.
\newblock \emph{ICML}, 2019.

\bibitem[Liu et~al.(2024)]{liu2024libero}
B.~Liu et~al.
\newblock LIBERO: Benchmarking Knowledge Transfer for Lifelong Robot Learning.
\newblock \emph{NeurIPS}, 2024.

\bibitem[Lu et~al.(2025)]{lu2025probing}
Y.~Lu et~al.
\newblock Linear Probing Reveals Task Structure in VLA Hidden Layers.
\newblock 2025.

\bibitem[Molinari et~al.(2025)]{molinari2025world}
M.~Molinari et~al.
\newblock Emergent World Representations in Vision-Language-Action Models.
\newblock 2025.

\bibitem[Oquab et~al.(2024)]{oquab2024dinov2}
M.~Oquab et~al.
\newblock DINOv2: Learning Robust Visual Features without Supervision.
\newblock \emph{TMLR}, 2024.

\bibitem[Radford et~al.(2021)]{radford2021clip}
A.~Radford et~al.
\newblock Learning Transferable Visual Models from Natural Language Supervision.
\newblock \emph{ICML}, 2021.

\bibitem[Rizzardo et~al.(2026)]{rizzardo2026observability}
C.~Rizzardo et~al.
\newblock Feature Observability and Controllability in VLAs.
\newblock 2026.

\bibitem[Swann et~al.(2026)]{swann2026sae}
E.~Swann et~al.
\newblock Sparse Autoencoder Interpretability for OpenVLA.
\newblock 2026.

\end{thebibliography}

\end{document}
